\subsubsection{Read the Original in Three Passes}

First identify the participants. \code{print-leaves} coordinates one request,
\code{tree2leaves} produces leaf replies, and \code{printer} performs output.
A future object buffers the replies between producer and consumer.
The file's first expression loads that library; the three object
definitions then establish the application's protocol.

Second read the consumer. Its \code{temporary} bindings belong to one
handler invocation: \code{leaf} holds the current value and \code{leaves}
is a fresh future. The dollar send starts the walker without waiting for
the entire traversal. Each \code{next-value} removes one reply or waits
until one arrives. The consumer forwards ordinary leaf values to the
printer and stops when it sees the keyword \code{:finish}.

Third read the producer's loop. Its \code{tree} variable is a working
representation of the remaining traversal, not an unchanged input tree.
The first two \code{match} clauses rewrite the remaining work.
The third sends one reply and removes the emitted head.
Only one matching clause executes on an iteration.

\begin{longtable}{L{0.27\textwidth}L{0.63\textwidth}}
\textbf{Pattern} & \textbf{How to read it}\\\hline
\endhead
\code{[[A] . B]} & The head is a one-element proper list.
Remove that extra layer, keeping the outer tail \code{B}.\\
\code{[[A . B] . C]} & The head is a nonempty pair.
Make its head and tail the next two pending items, followed by \code{C}.\\
\code{[A . B]} & The outer head is the next leaf.
Send \code{A} and continue with the outer tail \code{B}.\\
\end{longtable}

These patterns overlap. A one-element list is also a pair, so the first
clause must precede the second. A nested pair is also an outer head, so
the structural clauses precede the emission clause.
Uppercase names in the source are binding variables, not literal
symbols that the tree must contain.

\subsubsection{Trace a Small Tree Completely}

Use the data \code{'((a b) (c . d))}; it exercises both a proper nested
list and a dotted pair. The following is a logical trace of the walker,
not a claim about how its messages interleave with printing.
Unquoted uppercase letters in the table denote the already quoted
input symbols, not new variable lookups.

\begin{longtable}{L{0.30\textwidth}L{0.13\textwidth}L{0.32\textwidth}L{0.12\textwidth}}
\textbf{Before iteration} & \textbf{Clause} & \textbf{After iteration} &
\textbf{Reply}\\\hline
\endhead
\code{((A B) (C . D))} & 2 & \code{(A (B) (C . D))} & none\\
\code{(A (B) (C . D))} & 3 & \code{((B) (C . D))} & \code{A}\\
\code{((B) (C . D))} & 1 & \code{(B (C . D))} & none\\
\code{(B (C . D))} & 3 & \code{((C . D))} & \code{B}\\
\code{((C . D))} & 2 & \code{(C D)} & none\\
\code{(C D)} & 3 & \code{(D)} & \code{C}\\
\code{(D)} & 3 & \code{nil} & \code{D}\\
\end{longtable}

On the first iteration, pattern variables are \code{A = 'a},
\code{B = '(b)}, and \code{C = '((c . d))}.
The construction \code{[A B . C]} puts the value of \code{A} first,
the value of \code{B} second, and uses \code{C} as the final tail.
It therefore preserves the remaining work while exposing the next leaf.
When the working list becomes \code{nil}, the loop exits and the handler
sends one more reply: \code{:finish}.

The original sample's longer input adds further rewrites but follows
the same rules. It yields \code{A B C D E F G}.
No recursive call to \code{tree2leaves} is necessary: the loop and
rewrites implement the traversal.

\subsubsection{Where Does Completion Travel?}

Distinguish three message paths:
\begin{enumerate}[nosep]
  \item The top-level request goes to \code{print-leaves}. Its past-type
        send does not wait for a return value.
  \item Leaf values and the terminal keyword go from \code{tree2leaves}
        to the future's reply destination, then to the waiting consumer.
  \item The consumer sends each leaf to \code{printer} without requesting
        a reply. Printing is a side effect, not the future stream.
\end{enumerate}
The terminal marker tells the consumer that no more leaf values are part
of this request. It does not acknowledge completion of every printer
handler, and the consumer does not reply to its original sender.
Changing the outer invocation to a now-type send would therefore require
adding an application completion protocol; merely changing the arrow
would not create a reply.

\paragraph{Understanding check.}
Why is a leaf equal to the keyword \code{:finish} problematic?
The consumer would treat that data value as the terminal marker.
The original example assumes its data does not contain that marker.
Tagged values such as \code{[:value x]} solve the ambiguity, as in the
later stream exercise. Why not use \code{all-values} as a completion test?
Its buffer snapshot does not say whether the producer has finished.

\paragraph{A boundary of the original algorithm.}
The loop expects an outer proper list of pending items.
An empty nested list can be emitted as a \code{nil} leaf by the final
pattern; it is not automatically skipped as an empty subtree.
Read the actual patterns rather than assuming the behavior of a
general-purpose tree-flattening function.
